% !TeX root = ../../Book2.tex
\section{Fourier 变换}
\label{b2:sec:0301}

周期函数的频率由整数标记。若函数定义在整个欧氏空间上而不再周期重复，频率也变为连续变量。离散系数的求和于是让位于积分，但可积性、归一化与极限交换仍须逐项检查。

\subsection{定义与规范}
\label{b2:subsec:030101}

\begin{definition}\label{b2:def:03-fourier-transform}
设 \(d\) 为正整数，\(f\in L^1(\mathbb R^d;\mathbb C)\)。函数
\begin{equation}\label{b2:eq:03-fourier-definition}
 \widehat f(\xi)\coloneqq\mathcal Ff(\xi)
 =\int_{\mathbb R^d}f(x)\cdot e^{-2\pi\mathrm{i}x\cdot\xi}\,\mathrm{d}x
\end{equation}
称为 \(f\) 的\term[fourier-bianhuan]{Fourier 变换}[Fourier transform]。对 \(g\in L^1(\mathbb R^d)\)，另记
\[
 \check g(x)\coloneqq\int_{\mathbb R^d}g(\xi)\cdot e^{2\pi\mathrm{i}x\cdot\xi}\,\mathrm{d}\xi.
\]
此时只定义了正号积分；它何时反演 \(\mathcal F\)，将在\secref{b2:sec:0305}证明。
\end{definition}
\symidx{\(\mathcal Ff=\widehat f\)：欧氏空间 Fourier 变换}
\symidx{\(\check g\)：正号 Fourier 积分}

积分对每个 \(\xi\) 绝对收敛，且 \(|\widehat f(\xi)|\leq\|f\|_1\)，所以变换不依赖 \(f\) 的代表。这里 \(x\cdot\xi\) 是欧氏内积，虚数单位写作 \(\mathrm{i}\)。采用这个 \(2\pi\) 约定后，后面的反演公式没有额外常数，平方范数也将保持不变。若使用 \(e^{-\mathrm{i}x\cdot\xi}\) 的其他教材，需要同时转换频率尺度，不能只删去指数中的 \(2\pi\)。

\ProseParagraphBreak

在一维，取 \(a>0\)，区间示性函数给出一个可直接计算的模型：
\[
 \widehat{\mathbf1_{[-a,a]}}(\xi)=\dfrac{\sin(2\pi a\xi)}{\pi\xi}\quad(\xi\ne0),
 \quad \widehat{\mathbf1_{[-a,a]}}(0)=2a.
\]
原函数紧支集，并不使其变换紧支集；跳跃还使变换仅按 \(1/|\xi|\) 的尺度衰减。相反，对 \(a>0\)，分成正、负半轴积分可得
\[
 \mathcal F(e^{-a|x|})(\xi)=\dfrac{2a}{a^2+4\pi^2\xi^2}.
\]
这些公式同时提供了后面检验常数和反演条件的例子。

上面的连续频率可以用有限区间的级数作具体比较。在一维，设 \(f\) 连续且支撑在 \([-A,A]\)，并取 \(L>2A\)。将 \(f\) 在 \([-L/2,L/2]\) 上周期延拓，长度 \(L\) 的归一化 Fourier 系数为
\[
 c_L(k)=\frac1L\int_{-L/2}^{L/2}f(x)\cdot e^{-2\pi\mathrm i kx/L}\,\mathrm dx
       =\frac1L\widehat f(k/L).
\]
因此保留 \(|k|\leq RL\) 的有限频率和是
\[
 \frac1L\sum_{|k|\leq\lfloor RL\rfloor}
 \widehat f(k/L)\cdot e^{2\pi\mathrm i xk/L}.
\]
对固定 \(R>0\) 与 \(x\)，变换定义及控制收敛给出被求和函数在 \([-R,R]\) 连续；这一个和在 \(L\to\infty\) 时是频率步长 \(1/L\) 的 Riemann 和，趋于 \(\int_{-R}^R\widehat f(\xi)e^{2\pi\mathrm i x\xi}\,\mathrm d\xi\)。有限频带的极限只解释了积分为何自然出现；随后再令 \(R\to\infty\) 需要单独的收敛依据，不能从这个计算直接推出反演。归一化中的 \(1/L\) 正好成为频率积分元，所以采用 \(2\pi\) 后不再出现额外常数。

\subsection{平移}
\label{b2:subsec:030102}

对 \(f\in L^1(\mathbb R^d;\mathbb C)\)、\(h\in\mathbb R^d\)，沿用第一卷的平移约定 \(\tau_hf(x)\coloneqq f(x-h)\)。变量替换 \(y=x-h\) 给出
\[
 \widehat{\tau_hf}(\xi)=e^{-2\pi\mathrm{i}h\cdot\xi}\cdot\widehat f(\xi).
\]
平移保持 \(L^1\) 范数，故两侧都有定义。空间中的位置变化只改变频率中的相位，不改变 \(|\widehat f|\)。因而仅保留变换的模，通常不能恢复函数所在的位置。
\symidx{\(\tau_hf\)：平移，\(\tau_hf(x)=f(x-h)\)}

\subsection{调制}
\label{b2:subsec:030103}

对 \(\eta\in\mathbb R^d\)、\(f\in L^1(\mathbb R^d;\mathbb C)\)，算子
\[
 M_\eta f(x)\coloneqq e^{2\pi\mathrm{i}\eta\cdot x}\cdot f(x)
\]
称为\term[tiaozhi]{调制}[modulation]。合并指数便得
\[
 \widehat{M_\eta f}(\xi)=\widehat f(\xi-\eta).
\]
因此空间中的调制将频率整体平移。两种操作的次序要保留：\(\tau_hM_\eta=e^{-2\pi\mathrm{i}\eta\cdot h}M_\eta\tau_h\)，一般并不交换。
\symidx{\(M_\eta\)：调制算子}

\subsection{伸缩}
\label{b2:subsec:030104}

\begin{proposition}\label{b2:prop:03-linear-change}
设 \(d\) 为正整数，\(f\in L^1(\mathbb R^d;\mathbb C)\)，\(A\) 是可逆实 \(d\times d\) 矩阵，\(g(x)=f(Ax)\)，则
\[
 \widehat g(\xi)=|\det A|^{-1}\cdot\widehat f(A^{-\mathsf T}\xi).
\]
特别地，对 \(a>0\)，若 \(f_a(x)\coloneqq a^{-d}\cdot f(x/a)\)，则 \(\widehat{f_a}(\xi)=\widehat f(a\xi)\)。
\end{proposition}
\begin{proof}
在定义积分中令 \(y=Ax\)。体积元为 \(|\det A|^{-1}\mathrm{d}y\)，而 \(x\cdot\xi=y\cdot A^{-\mathsf T}\xi\)。绝对可积性由同一换元公式保证；取 \(A=a^{-1}I\) 再乘以归一化因子，即得后一式。
\end{proof}

一个在空间中压缩到宽度 \(a\) 的单位质量函数，其变换扩展到约 \(a^{-1}\) 的频率范围。这一换元公式精确记录了空间伸缩与频率伸缩的相反方向。


\begin{proposition}\label{b2:prop:03-symmetries}
设 \(d\) 为正整数，\(f\in L^1(\mathbb R^d;\mathbb C)\)，则 \(\widehat{\overline f}(\xi)=\overline{\widehat f(-\xi)}\)。若 \(f\) 实值且 \(f(-x)=f(x)\) 几乎处处，则其变换为实值偶函数。若对每个实正交矩阵 \(Q\) 都有 \(f(Qx)=f(x)\) 几乎处处，则 \(\widehat f(Q\xi)=\widehat f(\xi)\) 对每个 \(Q\) 及所有 \(\xi\in\mathbb R^d\) 成立。
\end{proposition}
\begin{proof}
第一式由对定义积分取共轭得到。偶性由反射换元得到，再与第一式组合即知实偶函数的变换为实值。对正交矩阵 \(Q\)，线性换元公式中的绝对行列式为 \(1\)、逆转置为 \(Q\)。若 \(f(Qx)=f(x)\)，便有 \(\widehat f(Q\xi)=\widehat f(\xi)\)。
\end{proof}

例如矩形 \(Q=\prod_j[-a_j,a_j]\) 的变换为各坐标区间变换的乘积，频率中的零点形成相应的超平面。球的示性函数则有旋转不变的变换；但旋转不变性并不提供其径向表达式，还须计算角积分。这里先由对称性确定变换的结构，具体特殊函数留待附录处理。

\ProseParagraphBreak

反射与共轭也是校验计算的工具。若一个实值偶函数被算出了非实的变换，或径向函数被算出了依赖坐标方向的变换，应先回查相位和换元。另一方面，实值函数不一定有实值变换：只有共轭对称是普遍结论，偶性是额外条件。

\subsection{微分与乘法}
\label{b2:subsec:030105}

\begin{proposition}\label{b2:prop:03-moments}
设 \(d\) 为正整数，\(m\) 为非负整数。若 \(f\in L^1(\mathbb R^d;\mathbb C)\) 且 \(|x|^m f\in L^1(\mathbb R^d)\)，则 \(\widehat f\in C^m(\mathbb R^d;\mathbb C)\)，并对 \(|\alpha|\leq m\) 有
\[
 \partial_\xi^\alpha\widehat f(\xi)
 =(-2\pi\mathrm{i})^{|\alpha|}\mathcal F(x^\alpha f)(\xi).
\]
单独考虑一个坐标 \(1\leq j\leq d\) 时，只要 \(f,x_jf\in L^1(\mathbb R^d;\mathbb C)\)，就有连续的偏导数 \(\partial_{\xi_j}\widehat f=-2\pi\mathrm{i}\mathcal F(x_jf)\)，无须假设其他坐标的一阶矩可积。
\end{proposition}
\begin{proof}
一阶差商中指数之差除以步长的模不超过 \(2\pi|x_j|\)。因此 \(|x_jf|\) 提供共同的可积控制，可以把差商极限移入积分。导数的连续性由相同控制下的控制收敛得到，这就证明了仅需一个坐标矩的结论。若具有 \(m\) 阶绝对矩，迭代时每次增加一个坐标因子；所有不超过 \(m\) 阶的单项式都由常数乘 \(1+|x|^m\) 控制，故每一步都合法。
\end{proof}

\begin{proposition}\label{b2:prop:03-derivative}
设 \(d\) 为正整数，\(1\leq j\leq d\)。若 \(f\in C^1(\mathbb R^d;\mathbb C)\)，且 \(f,\partial_jf\in L^1(\mathbb R^d)\)，则
\[
 \mathcal F(\partial_jf)(\xi)=2\pi\mathrm{i}\xi_j\widehat f(\xi).
\]
若 \(f\in L^1(\mathbb R^d;\mathbb C)\) 且其第 \(j\) 个弱导数属于 \(L^1(\mathbb R^d;\mathbb C)\)，公式也成立。
\end{proposition}
\begin{proof}
取光滑紧支集函数 \(\chi\)，在单位球上等于 \(1\)，并令 \(\chi_R(x)=\chi(x/R)\)。先对紧支集函数 \(\chi_Rf\) 分部积分，得到
\[
 \mathcal F(\chi_R\partial_jf)(\xi)
 +\mathcal F(f\partial_j\chi_R)(\xi)
 =2\pi\mathrm{i}\xi_j\mathcal F(\chi_Rf)(\xi).
\]
第一项及右侧用控制收敛取极限；中间误差的绝对值至多为 \(R^{-1}\|\partial_j\chi\|_\infty\cdot\|f\|_1\)，故趋零。弱导数情形直接把 \(\chi_Re^{-2\pi\mathrm{i}x\cdot\xi}\) 代入弱导数的测试恒等式，得到同一公式和估计。
\end{proof}

于是空间中的可积导数产生频率衰减，空间中的可积矩产生频率光滑性。两类假设各控制一个不同方向，不能把“变换连续”误认为原函数已经可微。
